\section{A microfounded interest-group example in the
Kaempfer--Lowenberg tradition}
\label{app:fourgroup}

The analytic example economies of \Cref{app:witness} are hand-chosen
parameter vectors. This section instead supplies a
public-choice interest-group polity in the \citet{Kaempfer1988} tradition,
with two interest groups in each country, in which the stage payoffs are
\emph{computed} from interest-group primitives by the
\citet{GrossmanHelpman1994} weighted-sum aggregation rule, and the
stage complementarity \eqref{eq:flip} and the diminishing-returns
curvature of \Cref{ass:concavity} then follow, on the sign region
defined below, from a \emph{sign condition on interest-group effectiveness}
rather than being posited directly at the stage level. The
construction is machine-checked in
\texttt{Theory.Microfoundations.FourGroupPolity} (the derivation and the
general-transition bridge \texttt{toParamsF}, which takes the transition law
and its five structural obligations as arguments) and instantiated in
\texttt{Examples.FourGroupWitness}.  The instance \texttt{toParams} under the partial-adjustment law --- the
same polity with its transition pinned to \texttt{affineF}~$\delta$ --- is the
\texttt{affineF}~$\delta$ instance of \texttt{toParamsF} and lives with the other affine
instances in \texttt{Examples.AffineInstance.Microfoundations}.

\subsection{Four groups, two per country}
\label{sec:fourgroup:primitives}

The sender and the target each contain two organized interest groups
--- four in all. On the sender side: a \emph{pro-sanction coalition}
$\JSgroup$ (strategic, security, or import-competing interests whose
political standing is advanced by confrontation) and a \emph{commercial
lobby} $\KSgroup$ (exporters and investors harmed by the sanction). On the
target side: a \emph{regime-aligned bloc} $\JTgroup$ that rallies under
external pressure, and an \emph{opposition} $\KTgroup$ that the sanction's
economic damage can empower. Each group $g$ has a welfare $W_g$ and a
political weight $\theta_g$ (its political effectiveness); following \citet{Kaempfer1988},
effectiveness is allowed to shift with the political state $x$ and the
sanction level $s$ as the confrontation hardens. Each country's
policymaker maximizes the \citet[eq.~(11)]{GrossmanHelpman1994}
weighted sum of its two groups' welfares, plus the resource terms that enter
its objective at unit weight --- on the sender side the sanctioning
deadweight $k_Ss^2$, the compliance--sanction cross term $w_{\mathrm{pl}}sa$ and the direct
compliance benefit $\beta^{\mathrm{base}}_Sa$; on the target side the
compliance cost $k_Ta^2$ and the sanction cost $d_ss$:
\begin{equation}
\Us=\!\!\sum_{g\in\{\JSgroup,\KSgroup\}}\!\!\theta_g\,W_g
+\beta^{\mathrm{base}}_S\,a-k_Ss^2-w_{\mathrm{pl}}\,s\,a,
\qquad
\Ut=\!\!\sum_{g\in\{\JTgroup,\KTgroup\}}\!\!\theta_g\,W_g
-k_Ta^2-d_s\,s ,
\label{eq:GH-aggregate}
\end{equation}
with group welfares $w_{\JSgroup}=g_Js$, $w_{\KSgroup}=-q_Ks$,
$w_{\JTgroup}=b_Jx-c_Ja$ and $w_{\KTgroup}=-b_Kx+c_Ka$. The unweighted terms
are the ones the planner of \Cref{sec:fourgroup:planner} keeps at unit
weight; they are not political transfers.

\subsection{Reduced stage payoffs}
\label{sec:fourgroup:payoffs}

Taking the linear-quadratic specification of the group welfares and
effectiveness schedules --- the pro-sanction coalition's pull rising
in the political state at rate $\nu$ (the source of the sender's
\emph{state--sanction complementarity},
$\theta_{\JSgroup}=\theta^0_{\JSgroup}+\nu x$), the regime bloc's and
opposition's effectiveness rising in the sanction at rates $\rho$ and
$\mu$ respectively, a sanctioning-cost curvature $k_S$, a compliance
cost $k_T$, and a cross-group \emph{compliance--sanction} term $w_{\mathrm{pl}}$
by which target compliance dampens the marginal political value of
sanctioning --- the aggregation \eqref{eq:GH-aggregate} collapses to
\begin{align}
\Us(x,a,s)&=\beta^{\mathrm{base}}_S\,a
+\bigl(\theta^0_{\JSgroup}+\nu x\bigr)g_J\,s
-\theta^0_{\KSgroup}q_K\,s
-k_S s^2
-w_{\mathrm{pl}}\,s\,a,
\label{eq:fourgroup-US}\\
\Ut(x,a,s)&=\bigl(\theta^0_{\JTgroup}+\rho s\bigr)\bigl(b_J x-c_J a\bigr)
+\bigl(\theta^0_{\KTgroup}+\mu s\bigr)\bigl(-b_K x+c_K a\bigr)
-k_T a^2-d_s s .
\label{eq:fourgroup-UT}
\end{align}
Every coefficient in \eqref{eq:fourgroup-US}--\eqref{eq:fourgroup-UT}
is an interest-group primitive; nothing is chosen to make a
monotonicity hold.

\subsection{The stage assumptions follow from a sign condition on interest-group effectiveness}
\label{sec:fourgroup:derived}

The cross-partials of \eqref{eq:fourgroup-US}--\eqref{eq:fourgroup-UT}
are read off directly:
\begin{equation}
\frac{\partial^2\Us}{\partial s\,\partial x}=\nu\,g_J,
\qquad
\frac{\partial^2\Us}{\partial s\,\partial a}=-w_{\mathrm{pl}},
\qquad
\frac{\partial^2\Us}{\partial s^2}=-2k_S,
\qquad
\frac{\partial^2\Ut}{\partial s\,\partial x}=\rho\,b_J-\mu\,b_K .
\label{eq:fourgroup-crosspartials}
\end{equation}
Define the \emph{sign region} $\mathcal R$ by the strict
sign conditions
\begin{equation}
k_S>0,
\qquad
\nu\,g_J>0,
\qquad
\mu\,b_K<\rho\,b_J,
\qquad
w_{\mathrm{pl}}\ge 0,
\label{eq:fourgroup-region}
\end{equation}
expressed entirely as inequalities among existing group primitives
(no new named gap scalar). On $\mathcal R$:
\begin{itemize}[leftmargin=2em]
\item \textbf{Stage complementarity (\Cref{ass:flip}).}
On $\mathcal R$, under the sign condition,
$\partial^2\Us/\partial s\,\partial x=\nu g_J>0$ is the sender's half of
the stage complementarity, and
$\partial^2\Ut/\partial s\,\partial x=\rho b_J-\mu b_K>0$ says that the
regime bloc's gain in effectiveness under sanctions strictly dominates opposition empowerment.
Both halves of \eqref{eq:flip} hold strictly --- this is exactly the
case-split \citeauthor{Kaempfer1988} leave with indeterminate sign,
here resolved by the region $\mathcal R$. (In Lean:
\path{U_S_senderStageIDStrict}, \path{U_T_targetStageIDStrict},
the derivative-free Topkis form that \eqref{eq:flip} implies.)
\item \textbf{\Cref{ass:concavity} (strict concavity of $\Us$ in $s$)
holds on $\mathcal R$ under the sign condition.} $\partial^2\Us/\partial s^2=-2k_S$, so the curvature
floor is $c_\sigma=2k_S>0$ whenever $k_S>0$. (In Lean:
\texttt{U\_S\_senderChordConcaveInS}; the same $k_S>0$ makes the
clipped-vertex static sender selector a genuine stage argmax,
\texttt{sigmaSel\_isMaxOn} --- the Rule-1 optimality coupling.)
\end{itemize}
\Cref{ass:flip} is thus a \emph{theorem} on $\mathcal R$, not a posit,
and the sign of the compliance--sanction cross term
$\partial^2\Us/\partial s\,\partial a=-w_{\mathrm{pl}}\le 0$ supplies
the decreasing-differences input the sender response uses.

The welfare identities and sign conditions of \Cref{sec:fourgroup:planner}
need none of this: they hold for every vector of the eighteen coefficients,
inside $\mathcal R$ or not (\path{FourGroupCoeffs}), and the planner-optimum
theorem there is stated for interest-group polities but its proof reads, of
$\mathcal R$, only $k_S>0$ and $w_{\mathrm{pl}}\ge0$. The region is what the
government game's stage assumptions need, so it still matters that $\mathcal R$
is inhabited without exhibiting a calibration, and it is. Of the eighteen
group-effectiveness coefficients, twelve --- the baseline weights
$\theta^0_{J_S},\theta^0_{K_S},\theta^0_{J_T},\theta^0_{K_T}$, the
anti-sanction coefficient $q_K$, the direct compliance benefit
$\mathrm{base}_S$, the concession costs $c_J,c_K$, the opposition weight
$b_K$, the opposition-empowerment slope $\mu$, the target concavity $k_T$
and the sanction-cost rate $d_s$ --- may be given \emph{arbitrary} real
values, and the remaining six are then obtained by choosing each strictly
beyond a threshold assembled from those already fixed: $k_S$ and
$w_{\mathrm{pl}}$ above zero, the factors $\nu$ and $g_J$ of the sender's
state--sanction cross term above zero, the hawk-welfare weight $b_J$ above zero, and finally
the slope $\rho$ of the regime bloc's effectiveness above $\mu b_K/b_J$, which is exactly the third
sign condition in \eqref{eq:fourgroup-region}. No
magnitude is selected anywhere in the construction (in Lean:
\path{exists_fourGroup_of_free}, with \path{Nonempty FourGroup} as a
one-line corollary). $\mathcal R$ is thus a sign region on interpretable
primitives rather than a condition whose satisfiability rests on one
parameter vector --- which is the appropriate standard for a hypothesis
class over which universally quantified results range.

It is worth being exact about what stage complementarity does \emph{not}
buy on its own. \Cref{prop:tipping} derives the tipping pivots
$\bar x_T,\bar x_S$ from \Cref{ass:flip} \emph{together with} the
boundary sign conditions \eqref{eq:boundary}, and
\Cref{ass:target} follows from those pivots, in the derivation after
\Cref{prop:flip-route}, only if the top sign $1\in\mathcal F_T(\bar s)$ holds --- the target's stage
payoff must already be non-decreasing in the sanction at the hardest
political state. The conditions \eqref{eq:fourgroup-region} that define
$\mathcal R$ restrict cross-partials and the sanction curvature; they
leave the \emph{level} of the target's sanction slope --- in particular
the sanction-cost rate $d_s$ --- entirely free, and so do not deliver
\eqref{eq:boundary}. On $\mathcal R$ the pivots are therefore derived
comparative statics wherever the boundary signs independently hold, and
are unavailable where they do not.

They do not hold at the calibrated certificate economy
\texttt{fgcParams} of \Cref{sec:fourgroup:instance}, and this is
machine-checked rather than conjectured. There the induced target stage
payoff is $\Ut(x,a,s)=3x-\tfrac15 a+s\,\bigl(x-\tfrac{14}{5}\bigr)$:
the target's state--sanction cross-partial $\rho b_J-\mu b_K=1>0$ is genuine ---
\Cref{ass:flip} does hold --- but the sanction-cost level
$d_s=\tfrac{14}{5}$ swamps it, so the sanction slope
$x-\tfrac{14}{5}\le-\tfrac95$ is strictly negative at every state and
action. Every chord in \eqref{eq:FT-def} therefore falls, and in
particular the top sign fails at every non-degenerate floor $\bar s<1$
($1\notin\mathcal F_T(\bar s)$:
\texttt{fgcCert\_one\_not\_mem\_targetFlipSet}), while the degenerate
floor $\bar s=1$ leaves only the trivial chord and so puts $0$ in the
flip set, destroying the bottom sign $0\notin\mathcal F_T$ instead;
either way the boundary bundle \eqref{eq:boundary} is contradictory
(\texttt{fgcCert\_stageFlip\_MS\_boundary\_infeasible}). The derived-basin
route therefore does \emph{not} fire at this economy: neither the
pivots of \Cref{prop:tipping} nor the derivation of \Cref{ass:target} along that
route is available here. The route is not thereby vacuous
--- it is inhabited by a separate economy designed for it (the
non-vacuity paragraph of \Cref{sec:flip}) --- but the paper refutes it
at \emph{this} economy rather than quietly assuming it.

\Cref{ass:target} at \texttt{fgcParams} is instead supplied by the
mutually exclusive \emph{entrenchment-condition} route, and supplied from the same
derived group primitives: the target's strictly negative stage slope is
bounded by the loss rate $c_T=\tfrac{14}{5}>0$ (numerically the derived
sanction cost $d_s$) and is offset by the continuation gain
$g_T=\tfrac{80}{11}$ carried on the target's invariant set of value functions, the entrenchment condition
$c_T=\tfrac{14}{5}\le\tfrac{36}{11}=\beta_T g_T\delta$ being discharged
outright (Lean: the fields \texttt{cT}, \texttt{gT},
\texttt{hT\_loss\_margin} of \texttt{fourGroupCertCertificate};
\Cref{sec:fourgroup:instance}). Stage complementarity and the sign
of the compliance--sanction cross term continue to do their work --- on the sender side, in the
comparative statics, and in the interest-group interpretation of the mechanism
--- but what they do not do at this calibration is generate the basin,
which the entrenchment condition generates instead. Finally, what is
\emph{not} discharged unconditionally from \eqref{eq:fourgroup-region}
alone is the \emph{dynamic} content of \Cref{ass:sender} and
\Cref{ass:target} --- single-crossing of the forward-looking objectives
once the continuation value is included --- together with joint MPE
existence; these are supplied on an invariant set of value functions by a
model-specific row-Lipschitz / Schur-condition route
(\Cref{sec:existence,app:existencemachinery}).

\subsection{A concrete instance and its long-run behavior}
\label{sec:fourgroup:instance}

The Lean parameter vector \texttt{fourGroupWitnessParams} fixes
\begin{equation}
\nu g_J=3,
\quad
\rho b_J-\mu b_K=\tfrac{9}{10},
\quad
w_{\mathrm{pl}}=\tfrac{3}{2},
\quad
c_\sigma=2k_S=1,
\quad
\delta=\tfrac15,\ \beta_T=\tfrac{9}{10},\ \beta_S=\tfrac12,
\end{equation}
all four region conditions \eqref{eq:fourgroup-region} holding
strictly. The
machine-checked headline
\path{fourGroupWitnessParams_three_goals_conditional} delivers
the dynamic-over-myopic comparison on the basin, conditional on the value class, with
the inputs of the derived-basin route feeding it discharged from
\eqref{eq:fourgroup-crosspartials}.

\paragraph{The conditional headline is now discharged unconditionally.}
The value-class residual that
\texttt{fourGroupWitnessParams\_three\_goals\_conditional} leaves open
is closed by a certificate-complete sibling, another interest-group
polity. At the tuned instance \texttt{fgcParams} --- the
\citet{GrossmanHelpman1994} interest-group polity with $\delta=\tfrac12$,
$\beta_T=\tfrac{9}{10}$, $\beta_S=\tfrac12$ and the region conditions
holding strictly --- every field of the core certificate
\texttt{StandardComplianceLocalCoreCertificate} is discharged from the
interest-group primitives: the sender high-sanction condition with
$\beta_S K_S\delta=\tfrac13<\tfrac32=m_S$, and the target entrenchment condition
with $c_T=\tfrac{14}{5}\le\tfrac{36}{11}=\beta_T g_T\delta$, on the
slope-pinned invariant sets
\begin{align*}
\mathcal C_S&=\{V:\lVert V\rVert_\infty\le 5,\ V(y)-V(x)=\tfrac43\,(y-x)\ \text{for } 0\le x\le y\le1\},\\
\mathcal C_T&=\{V:\lVert V\rVert_\infty\le 28,\ V(y)-V(x)=\tfrac{80}{11}\,(y-x)\ \text{for } 0\le x\le y\le1\}
\end{align*}
(each is a Lipschitz ball whose lower slope floor on the whole physical
interval equals its Lipschitz ceiling, so the set is pinned to one slope;
the radii are the discounted stage bounds $\tfrac{5/2}{1-1/2}=5$ and
$\tfrac{14/5}{1-9/10}=28$; \path{fgcCarrierS_eq}, \path{fgcCarrierT_eq}).
The resulting headline
\texttt{fourGroupCert\_three\_goals\_dominance} (axiom-clean) is
\emph{unconditional}: it \emph{constructs} the stationary feedback
Stackelberg equilibrium $E$ on those invariant sets and delivers the
dynamic-over-myopic comparison (each side at its own equilibrium's target
action) $\sigM(x,\alphaM(x))\le\sigDynTop_{\Vs}(x,\alpha(x))$ for every
$x\in[\tfrac1{10},1]$, with no value-class certificate taken as a
hypothesis. This is the first \emph{microfounded} instance of the core
certificate --- existence and the three goals hold at a
public-choice-derived stage economy, not only at the analytic
example economies of \Cref{sec:existence:nonvacuity}.

\paragraph{What the certificate-complete economy exhibits, and what it does not.}
At \texttt{fgcParams} the two Bellman fixed points are exact and affine,
$\Vs(x)=\tfrac{11}{3}+\tfrac43x$ and $\Vt(x)=\tfrac{52}{11}+\tfrac{80}{11}x$
(\path{fgc_VS_star_fixed}, \path{fgc_VT_star_fixed}), and the derived
actions of the constructed equilibrium are full non-compliance and the maximal
sanction at every physical state (\path{fgc_}$\alpha$\path{_on_eq_zero},
\path{fgc_}$\sigma$\path{_dyn_zero_eq_one}: the joint objective, with the
sender's dynamic best response substituted, is maximized at full non-compliance, and at full
non-compliance the sender's dynamic best response is $1$). Its dynamic-over-myopic
comparison is an equilibrium comparison of the constructed equilibrium with the
myopic equilibrium, in which both governments re-optimize; it is not the
sender's foresight premium at a common compliance level, since the myopic sanction
$\min(1,2+x)$ and the dynamic best response at full non-compliance are both $1$. This economy supplies
microfounded non-vacuity of the core certificate and the welfare comparison
of \Cref{sec:fourgroup:planner}. Whether either government's
forward-looking incentive changes its behavior at a given state is a
separate question that this example economy is not used to answer; the uniform
statement that both forward-looking incentives change behavior belongs to the class under the lifting condition
(\Cref{prop:placable-family} and its form without the one-period ranking condition after
\Cref{thm:four-cells}). The earlier calibration
\texttt{fourGroupWitnessParams} of \Cref{sec:fourgroup:instance}, whose
bistability was explored numerically only, supports no machine-checked
statement about its continuation values or the size of any foresight premium, and none
is made here. Displaying the indirect effect of the target's lower compliance at an interior target policy
remains the role of the mechanism-active $P_3$ economy (\Cref{sec:existence:nonvacuity}).

\paragraph{What this example adds.}
The interest-group example \emph{derives} the two stage-level primitives ---
\Cref{ass:flip} (stage complementarity) and \Cref{ass:concavity} (strict
sender concavity in the sanction) --- from public-choice primitives, so the
``flip'' in the cross-partials is a consequence of the regime bloc's gain in effectiveness under
sanctions exceeding opposition empowerment rather than an analyst's posit. What it does
\emph{not} do is obtain \Cref{ass:target} from \Cref{ass:flip}: the
derived-basin route of \Cref{prop:flip-route} needs the boundary sign conditions
\eqref{eq:boundary} \emph{in addition} to complementarity, and the region
conditions \eqref{eq:fourgroup-region} constrain cross-partials and
curvature, not the \emph{level} of the target's sanction slope
(\Cref{sec:fourgroup:derived}). At the calibration
\texttt{fourGroupWitnessParams} the target half of \Cref{ass:target} is
carried inside the value-class certificate that the conditional headline
assumes; at the certificate-complete \texttt{fgcParams} it is discharged
outright on the mutually exclusive entrenchment-condition route ($c_T=\tfrac{14}{5}>0$),
the boundary bundle of the derived-basin route being machine-checked infeasible there. That
instance also discharges dynamic single-crossing and joint MPE existence on
its invariant set of value functions and \emph{constructs} the equilibrium
(\texttt{fourGroupCert\_three\_goals\_dominance}). This is the decoupled
$L_{ST}=0$ route; a coupled three-goals/Schur certificate would need
analytic discharges of the conditions on the target's invariant set that are not available
on the sets examined, and none is claimed.

The grounding is a \emph{trade}, not a one-directional subsumption,
relative to the Eaton--Engers example economy of \Cref{sec:existence:nonvacuity}:
that economy derives the \emph{transition} sign from a two-group
Kaempfer--Lowenberg market-clearing polity (the assignment of states to the lock-in
basin and to the basin of convergence to zero is an equilibrium output, not the bare partial-adjustment law),
whereas the interest-group example derives the \emph{stage payoffs} but
takes the partial-adjustment transition as given. The two constructions ground
complementary halves of the model; neither subsumes the other.

\subsection{The pressure schedules from a Becker contest}
\label{sec:fourgroup:contest}

The market-clearing polity above takes the groups' pressure schedules
as monotone primitives. They too can be derived. Let the two
target-side groups choose pressure levels at quadratic cost, with
political influence additive in the pressure difference --- the
canonical pressure-competition structure of \citet{Becker1983}. Best
responses are then \emph{dominant strategies} with closed-form
equilibrium pressure $p_i^\ast = \max\{0,\,S_i(s)/(2c_{p_i})\}$, where
the stake $S_i(s)$ is the group's effectiveness-weighted marginal
benefit of entrenchment net of income-incidence erosion --- both read
off the interest-group payoff fields already introduced. The three
Kaempfer--Lowenberg effects (rally, incidence, empowerment) become
\emph{definitions} --- equilibrium contest objects rather than assumed
schedules --- and the market-clearing map they induce is exactly the
partial-adjustment transition used throughout, so the lock-in
condition becomes a \emph{cost-adjusted} form of the third sign condition of region~(R),
$c_{p_J}(\mu b_K) < c_{p_K}(\rho b_J)$, degenerating literally to that
sign condition at equal pressure costs. The
only new primitives are the two pressure-cost coefficients and two
incidence-exposure slopes; the lock-in conclusion is obtained by
\emph{applying} the existing market-clearing theorem, and at the
example's numbers the contest-induced transition coincides with the
partial-adjustment law --- the transition used throughout the paper is, at the
example economy, itself a contest-equilibrium object.

\subsection[Aggregate welfare: construction of W and the two comparisons]
{Aggregate welfare: construction of $W$, the two
comparison paths, and the period-by-period accounting}
\label{sec:fourgroup:planner}

\Cref{sec:policy:welfare} states the welfare results on the class; this
subsection carries the polity's instance of them and the derivation behind it. Aggregate
welfare, its closed form at the example economy, the two state paths, the
per-period gap and every discounted constant below are machine-checked
(\path{Theory.Microfoundations.FourGroupPlanner},
\path{Examples.FourGroupPlannerWitness} and, for the primitive identity,
its sign conditions and the two counterexamples,
\path{Examples.FourGroupWelfareSign}, and, for the planner's optimum at every
polity and transition, \path{Theory.Microfoundations.FourGroupPlannerOptimum};
\Cref{app:formalization}). The accounting identities, the sign conditions and
the stationary comparison are stated over the coefficients alone, with no
region condition (\path{FourGroupCoeffs}), and a polity reaches them through
\path{FourGroup.toCoeffs}; the
one-line stationarity reduction $x=\F(x,s)\iff x=s$ is immediate from
\eqref{eq:affine}.

\paragraph{Construction of aggregate welfare.}
The \citet{GrossmanHelpman1994} rule \eqref{eq:GH-aggregate} weights each
group's welfare by that group's political effectiveness $\theta_g$. In the
reduced payoffs \eqref{eq:fourgroup-US}--\eqref{eq:fourgroup-UT} those
weights appear not as bare constants but as the \emph{schedules}
$\theta_{\JSgroup}(x)=\theta^0_{\JSgroup}+\nu x$,
$\theta_{\KSgroup}=\theta^0_{\KSgroup}$,
$\theta_{\JTgroup}(s)=\theta^0_{\JTgroup}+\rho s$ and
$\theta_{\KTgroup}(s)=\theta^0_{\KTgroup}+\mu s$. Aggregate welfare is the \emph{same}
aggregation with the political weights removed: each schedule --- not
merely its intercept --- set identically to one, so the effectiveness
slopes $\nu,\rho,\mu$ vanish along with them. Nothing else
changes: the group-borne resource terms already enter
\eqref{eq:fourgroup-US}--\eqref{eq:fourgroup-UT} at unit weight,
multiplied by no $\theta_g$, so they carry through untouched. Substituting
$\theta_{\JSgroup}\equiv\theta_{\KSgroup}\equiv
\theta_{\JTgroup}\equiv\theta_{\KTgroup}\equiv1$
into \eqref{eq:fourgroup-US}--\eqref{eq:fourgroup-UT} and summing the two
countries' group welfares gives
\begin{equation}
W(x,a,s)
=\bigl(g_J\,s-q_K\,s+\beta^{\mathrm{base}}_S\,a-k_S s^2
        -w_{\mathrm{pl}}\,s\,a\bigr)
+\bigl((b_J-b_K)\,x+(c_K-c_J)\,a-k_T a^2-d_s s\bigr),
\label{eq:planner-W}
\end{equation}
the first bracket the two sender groups' summed welfare, the second the
two target groups'. Every coefficient in \eqref{eq:planner-W} is an
existing group primitive of
\eqref{eq:fourgroup-US}--\eqref{eq:fourgroup-UT}; no scalar is introduced.
Subtracting \eqref{eq:planner-W} from the sum of \eqref{eq:fourgroup-US}
and \eqref{eq:fourgroup-UT} and collecting group by group returns
\eqref{eq:planner-distortion} coefficient by coefficient in all eighteen
fields --- pure algebra, no side condition. Evaluating
\eqref{eq:planner-W} at the two corners the dynamics single out gives
$W(0,0,0)=0$ and $W(1,0,1)=(g_J-q_K)+(b_J-b_K)-(k_S+d_s)$, and comparing
the two is the flip boundary reported in \Cref{sec:policy:welfare}.
For any planner discount factor $\beta\in[0,1)$, capitalizing a
stationary flow gives
\[
\mathcal W_\beta^{\mathrm{stat}}(x,a,s)
=\frac{W(x,a,s)}{1-\beta}.
\]
Therefore
\[
\mathcal W_\beta^{\mathrm{stat}}(0,0,0)
-\mathcal W_\beta^{\mathrm{stat}}(1,0,1)
=
\frac{(k_S+d_s)-\bigl[(g_J-q_K)+(b_J-b_K)\bigr]}
     {1-\beta},
\]
so the discounted stationary ranking has exactly the same parametric
boundary as the flow ranking. Lean proves the capitalization identity,
the discounted if-and-only-if boundary, and the exact primitive gap as
\path{stationaryPlannerValue_eq},
\path{stationaryPlannerValue_lockin_dominated_iff}, and
\path{stationaryPlannerValue_relief_gain}.

At the certificate-complete instance \texttt{fgcParams} of
\Cref{sec:fourgroup:instance} the sender groups' sanction benefits net out
($g_J=q_K=1$), the target groups' net entrenchment value is
$b_J-b_K=1$, the direct compliance flows cancel
($\beta^{\mathrm{base}}_S=0$, $c_J=c_K$) and compliance curvature is
absent ($k_T=0$), so \eqref{eq:planner-W} collapses to
\begin{equation}
W(x,a,s)=x-\tfrac{14}{5}s-\tfrac12 s^2-\tfrac{11}{5}\,s\,a :
\label{eq:planner-W-fgc}
\end{equation}
entrenchment is worth $x$, and sanctioning burns the target's sanction
cost $d_s=\tfrac{14}{5}$, the sender's deadweight $k_S s^2=\tfrac12 s^2$
and the compliance--sanction transfer $w_{\mathrm{pl}}\,s\,a=\tfrac{11}{5}sa$. On the
physical cube $[0,1]^3$, \eqref{eq:planner-W-fgc} obeys
$\lvert W\rvert\le\tfrac{13}{2}$ --- the uniform bound the discounting
step below consumes.

\paragraph{Stationary points of the transition.}
Under the partial-adjustment law \eqref{eq:affine} with $\delta\in(0,1)$,
$x=\F(x,s)$ holds iff $x=s$: the stationary points are exactly the
diagonal triples $(s,a,s)$, $(s,a)\in[0,1]^2$. Evaluating
\eqref{eq:planner-W-fgc} along that diagonal, the entrenchment credit and
the sanction cost collapse into a single linear term,
\[
W(s,a,s)=-\tfrac95 s-\tfrac12 s^2-\tfrac{11}{5}\,s\,a\;\le\;0=W(0,0,0),
\]
each of the three terms being nonpositive on the cube, with \emph{strict}
inequality whenever $s>0$ (the first is then strictly negative). So no
stationary point does better than the origin, and every stationary point
carrying \emph{positive sanctions} is strictly worse. That weak-with-
strict-above-zero statement is exactly what is claimed and no more: the
zero-sanction stationary points $(0,a,0)$ all attain $0$ as well, so what
is pinned is that \emph{sanctioning} at a stationary point is what
destroys planner value --- not that the origin is the unique maximizer.
The lock-in corner $(1,0,1)$ --- the absorbing point of the certified play
--- sits $\tfrac{23}{10}$ below the origin in flow terms.

\paragraph{The two comparison paths.}
At \texttt{fgcParams}, $\delta=\tfrac12$, so \eqref{eq:affine} reads
$\F(x,s)=(x+s)/2$. The \emph{equilibrium} path is the orbit of the
certified play, not of an assumed policy: the derived sender corner
$\sigDynTop_{\Vs}(x,0)=1$ against the derived leader policy
$\alpha\equiv 0$ (full non-compliance), both pinned for every value function in
the sender's certified invariant set and every state in $[0,1]$ by the certificate
behind \texttt{fourGroupCert\_three\_goals\_dominance}, so
$x_{t+1}=\F(x_t,1)=(x_t+1)/2$. The \emph{comparison} path shuts sanctions
off, $s\equiv0$, $a\equiv0$ --- a \emph{feasible} policy, not an
equilibrium one, which is all an inefficiency claim requires --- so
$x_{t+1}=\F(x_t,0)=x_t/2$. Both are affine contractions and solve in
closed form from any start $x_0$:
\begin{equation}
x^E_t=1-(1-x_0)\,2^{-t}\uparrow 1,
\qquad
x^0_t=x_0\,2^{-t}\downarrow 0 .
\label{eq:planner-paths}
\end{equation}
Both orbits stay in $[0,1]$. The cube-wide policy identities proved for the
exact equilibrium (\path{fgcEquilibriumTargetPath_eq_zero},
\path{fgcEquilibriumSenderPath_eq_one}) identify the equilibrium path with
equilibrium play from \emph{every} start $x_0\in[0,1]$; for a basin start
$x_0\in[\tfrac1{10},1]$ the orbit is increasing and also remains in the high
basin $[\tfrac1{10},1]$ used by the separate three-goals certificate, an
additional invariance fact and not the domain of the welfare theorem.

\paragraph{Period-by-period accounting and the discounted gap.}
Substituting \eqref{eq:planner-paths} into \eqref{eq:planner-W-fgc}, the
equilibrium flow is
$W(x^E_t,0,1)=x^E_t-\tfrac{33}{10}=-\tfrac{23}{10}-(1-x_0)2^{-t}$
(entrenchment worth $x^E_t$ against the per-period resource burn
$k_S+d_s=\tfrac{33}{10}$) and the comparison flow is
$W(x^0_t,0,0)=x_0 2^{-t}$, so the per-period gap is
\begin{equation}
W(x^0_t,0,0)-W(x^E_t,0,1)=\tfrac{23}{10}+2^{-t},
\label{eq:planner-gap}
\end{equation}
in which the start $x_0$ cancels \emph{identically}: a permanent flow gain
of $\tfrac{23}{10}$ from abandoning the sanctions war, plus a vanishing
transition bonus. Now discount at any $\beta\in[0,1)$ and sum. Along any
cube-valued path the series $\sum_t\beta^t W$ converges absolutely by
geometric comparison against $\tfrac{13}{2}\sum_t\beta^t$, so both path
values are well defined and may be summed term by term. Each is then a
pair of geometric series, with ratios $\beta$ (the constant flow) and
$\beta/2$ (the transition term inherited from the contraction factor
$1-\delta=\tfrac12$ in \eqref{eq:planner-paths}):
\begin{equation}
V^E(x_0;\beta)=-\frac{23}{10(1-\beta)}-\frac{1-x_0}{1-\beta/2},
\qquad
V^0(x_0;\beta)=\frac{x_0}{1-\beta/2}.
\label{eq:planner-values}
\end{equation}
Differencing \eqref{eq:planner-values}, the two $(1-\beta/2)^{-1}$ terms
recombine --- $x_0+(1-x_0)=1$ --- and the start drops out, leaving
\[
V^0(x_0;\beta)-V^E(x_0;\beta)
=\frac{23}{10(1-\beta)}+\frac{1}{1-\beta/2},
\]
the exact discounted gap of \Cref{sec:policy:welfare}, an identity for
every start $x_0\in[0,1]$ and every $\beta\in[0,1)$; both terms are
strictly positive on that range, which is the strict-dominance form.

\paragraph{The stationary comparison.}
Let $g_J$ and $q_K$ be the sender groups' per-unit sanction benefit and
loss, $b_J$ and $b_K$ the target groups' state valuations, $k_S$ the
sender's sanctioning-cost curvature, and $d_s$ the target's sanction-cost
rate. The benefit of maximal entrenchment and sanctions is
$(g_J-q_K)+(b_J-b_K)$; its resource cost is $k_S+d_s$. For a planner discount
factor $\beta\in[0,1)$, define discounted welfare at a constant
state-action-sanction triple by
$\mathcal W_\beta^{\mathrm{stat}}(x,a,s)=\sum_{t=0}^{\infty}\beta^t W(x,a,s)$.

\begin{proposition}[Discounted stationary welfare boundary]
\label{prop:discounted-stationary-welfare}
For every vector of interest-group coefficients, inside region~(R) or not, and
every $\beta\in[0,1)$,
\[
\mathcal W_\beta^{\mathrm{stat}}(0,0,0)
-\mathcal W_\beta^{\mathrm{stat}}(1,0,1)
=
\frac{(k_S+d_s)-\bigl[(g_J-q_K)+(b_J-b_K)\bigr]}
     {1-\beta}.
\]
Consequently, the stationary zero path strictly planner-dominates the
stationary lock-in path if and only if
\[
(g_J-q_K)+(b_J-b_K)<k_S+d_s.
\]
\end{proposition}

Discounting capitalizes the flow gap without changing its sign, and the
example interest-group economy satisfies the inequality. No region condition is
read (\path{FourGroupCoeffs.stationaryPlannerValue_relief_gain},
\path{FourGroupCoeffs.stationaryPlannerValue_lockin_dominated_iff}). The stationary
comparison is not the equilibrium comparison, which runs along moving paths
from a common start; that comparison is the next proposition. (This
proposition stood in the main text until round AD; it was moved here, with
the affine closed form, when the welfare section was reduced to one attained
example and the family-wide statement.)

\paragraph{The same accounting in the primitives, and the sign.}
The two closed forms are the case $\delta=\tfrac12$, $b_J-b_K=1$,
$(k_S+d_s)-(g_J-q_K)=\tfrac{33}{10}$ of an identity that holds for every
vector of the eighteen coefficients and every transition of the paper's class.
No sign condition of region~(R) enters its proof, so it also covers
coefficient vectors outside the region, such as those in which opposition
empowerment dominates the regime bloc's gain in effectiveness under sanctions (\path{FourGroupCoeffs}). Write $H=(k_S+d_s)-(g_J-q_K)-(b_J-b_K)$ for the
stationary gap and $B=b_J-b_K$ for the target groups' net valuation of the
state.

\begin{proposition}[The welfare gap in primitives]
\label{prop:affine-welfare-gap}
For every vector of interest-group coefficients, every transition
$\F\colon[0,1]^2\to[0,1]$ satisfying \textup{(F2)--(F3)}, every planner
discount $\beta\in[0,1)$ and every start $x_0\in[0,1]$, let
$e_{t+1}=\F(e_t,1)$ and $y_{t+1}=\F(y_t,0)$, with $e_0=y_0=x_0$, be the
corner path of full sanctions against full non-compliance and the zero-sanction path, and let $z_t=1-e_t+y_t$, so
that $z_t\in[0,1]$ and $z_0=1$. Then
\[
V^0(x_0;\beta)-V^E(x_0;\beta)
=\frac{H}{1-\beta}+B\sum_{t=0}^{\infty}\beta^tz_t .
\]
If the transition also satisfies \textup{(F1)}, \textup{(F4)} and strict
adjustment \textup{(F5)}, the difference is positive for every
$\beta\in[0,1)$ if and only if $k_S+d_s>g_J-q_K$ and
$k_S+d_s\ge(g_J-q_K)+(b_J-b_K)$. Under the partial-adjustment law \eqref{eq:affine} with
$\delta\in(0,1)$ the two paths are $x^0_t=x_0(1-\delta)^t$ and
$x^E_t=1-(1-x_0)(1-\delta)^t$, the residual is $z_t=(1-\delta)^t$, and
\[
V^0(x_0;\beta)-V^E(x_0;\beta)
=\frac{(k_S+d_s)-(g_J-q_K)-(b_J-b_K)}{1-\beta}
+\frac{b_J-b_K}{1-\beta(1-\delta)} ;
\]
this difference is positive if and only if
$[(k_S+d_s)-(g_J-q_K)-(b_J-b_K)](1-\beta(1-\delta))+(b_J-b_K)(1-\beta)>0$,
and it is positive for every $\beta\in[0,1)$ if and only if
$k_S+d_s>g_J-q_K$ and $k_S+d_s\ge(g_J-q_K)+(b_J-b_K)$.
\end{proposition}

\begin{proof}
At full non-compliance $W(x,0,s)=(g_J-q_K-d_s)s-k_Ss^2+(b_J-b_K)x$, so the
per-period gap between the two paths is $H+Bz_t$. Both paths stay in $[0,1]$,
where the flows are bounded, and summing termwise gives the identity
(\path{FourGroupCoeffs.pathValue_cornerF_gap}). Under strict adjustment the
corner path rises to $1$ and the zero-sanction path falls to $0$, so their
separation $1-z_t$ starts at $0$, stays in $[0,1]$ and tends to one; with
$C=(k_S+d_s)-(g_J-q_K)$, a discounted series $\sum_t\beta^t\bigl(C-B(1-z_t)\bigr)$
of that kind is positive at every $\beta\in[0,1)$ exactly when $C>0$ and
$B\le C$ (\path{DiscountedSign.tsum_pos_all_discounts_iff},
\path{pathValue_cornerF_gap_pos_all_discounts_iff_of_params}). Under the partial-adjustment
law the residual is $(1-\delta)^t$ and the start cancels; two geometric series
give the closed form (\path{pathValue_affineCorner_gap}). Clearing the
positive denominators gives the fixed-$\beta$ condition
(\path{affineCornerGap_pos_iff}). At $\beta=0$ the numerator is
$(k_S+d_s)-(g_J-q_K)$, and as $\beta\to1$ it tends to
$\delta[(k_S+d_s)-(g_J-q_K)-(b_J-b_K)]$, which gives the all-$\beta$ condition
(\path{affineCornerGap_pos_all_discounts_iff}). The welfare theorems cited are
stated over the coefficients and read no region condition; the polity forms are
their instances through \path{FourGroup.toCoeffs}.
\end{proof}

The conjunction ``stationary gap positive and $b_J-b_K\ge0$'' is sufficient
and not necessary: a zero stationary gap with a positive transition bonus,
or a negative $b_J-b_K$, can still give a positive difference. Since
$1\le\sum_t\beta^tz_t\le1/(1-\beta)$, the difference is at least
$H/(1-\beta)+B$ when $B\ge0$ and at least $[(k_S+d_s)-(g_J-q_K)]/(1-\beta)$
when $B<0$ (\path{pathValue_cornerF_gap_ge_of_nonneg},
\path{pathValue_cornerF_gap_ge_of_neg}). A lower bound on the transition's sanction slope bounds the
residual sum from above and so gives no lower bound in the first case; a lower
bound $g\ge0$ on the transition's state slope along the two endpoint sections sharpens the first bound's
transition term to $B/(1-\beta g)$
(\path{pathValue_cornerF_gap_ge_of_nonneg_of_cornerChords}). The all-discount
condition needs no affine rate: it holds under the qualitative direction of
\textup{(F5)} with the continuity of \eqref{eq:transition}, which is the
paper's maintained class, and under its quantitative form, a positive lower
bound on the sanction slope, such as the family's $d_{\mathrm{lo}}>0$, which adds a
geometric rate (\path{pathValue_cornerF_gap_pos_all_discounts_iff}). Without
strict adjustment the necessity direction fails: the identity transition keeps
both paths at the start and gives $[(k_S+d_s)-(g_J-q_K)]/(1-\beta)$ whatever
$b_J-b_K$ is (\path{not_necessary_without_strict_adjustment}).

\paragraph{Region~(R) does not determine the sign.}
Neither inequality of \Cref{prop:affine-welfare-gap} is a region-(R)
condition, and the government game does not settle them either. Rescale the
pro-sanction coalition: multiply its per-unit benefit $g_J$ by some $\phi\ne0$
and divide its weight schedule, $\theta^0_{\JSgroup}$ and $\nu$, by the same
$\phi$ (\path{hawkRescale}). The product $(\theta^0_{\JSgroup}+\nu x)g_J$ does not
change, so both government payoffs, the static sender selector, every
region-(R) condition and the induced game under every transition are those of
the original polity, and every equilibrium of the one is an equilibrium of the
other (\path{hawkRescale_U_S}, \path{hawkRescale_U_T},
\path{hawkRescale_toParamsF}). The planner, who counts the hawks' benefit at
unit weight, sees $W+(\phi-1)g_Js$, so under every transition into $[0,1]$ the
gap between the zero-sanction path and the corner path of full sanctions
against full non-compliance falls by
exactly $(\phi-1)g_J/(1-\beta)$ (\path{hawkRescale_cornerF_gap}), while under
\textup{(F2)--(F3)} the gap is at most $(|H|+|B|)/(1-\beta)$ in absolute value
(\path{abs_cornerF_gap_le}). Hence, for every polity whose hawks gain from
sanctions, $g_J>0$, some $\phi>0$ makes the gap strictly negative at every planner
discount, every start and every transition satisfying \textup{(F2)--(F3)}
(\path{exists_hawkRescale_cornerF_gap_neg}); for every polity whose hawks lose,
$g_J<0$, which region~(R) allows when $\nu<0$, some $\phi>0$ makes it strictly
positive (\path{exists_hawkRescale_cornerF_gap_pos}). Each statement goes one
way: when the hawks gain, the game never shows that the zero-sanction path
is at least as good for the planner as persistent full sanctions, and when
they lose, it never shows the reverse.

Next to the example economy, the rescaling with $\phi=5$ raises $g_J$ from $1$
to $5$ and scales the weight schedule down to $\theta^0_{\JSgroup}=\tfrac35$,
$\nu=\tfrac15$, so that $(\tfrac35+\tfrac x5)\cdot5=3+x$ is exactly what it was
(\path{fourGroupCertHawkScaled_eq_hawkRescale}). The induced game is
\texttt{fgcParams} itself and the constructed equilibrium is an equilibrium of
the new polity verbatim (\path{fourGroupCertHawkScaled_toParams}), while the
planner sees $W+4s$: $(k_S+d_s)-(g_J-q_K)=-\tfrac7{10}$, the gap falls by
$4/(1-\beta)$ (\path{fourGroupCertHawkScaled_gap_shift}), and the value of the
specified feasible zero-sanction path less that of equilibrium play is
\[
-\frac{17}{10(1-\beta)}+\frac{1}{1-\beta/2}
=-\frac{14+3\beta}{20(1-\beta)(1-\beta/2)}<0
\]
at every $\beta\in[0,1)$ (\path{fourGroupCertHawkScaled_gap_neg}): same
government game, same equilibrium, opposite welfare ranking. The reversal
concerns the same specified zero-sanction alternative, not the hawk-scaled
planner's own optimum. Bundled, the two polities induce \texttt{fgcParams}
and their planners, who maximize aggregate welfare, rank the equilibrium path and the same
feasible zero-sanction alternative oppositely at every discount and every
physical start
(\path{fourGroupCert_welfare_not_identified_by_game}): government payoff
functions and equilibrium behavior do not by themselves determine the
welfare ranking of these two paths. A second polity
changes only the direct sanction cost to $d_s=\tfrac1{10}$, so region~(R)
again holds and the cost stays positive: the difference is $\tfrac35$ at
$\beta=0$, zero at $\beta=\tfrac34$ and $-\tfrac{24}{11}$ at
$\beta=\tfrac9{10}$ (\path{fourGroupCertPositiveCost_gap_at_nine_tenths});
the difference equals $(3-4\beta)/(5(1-\beta)(1-\beta/2))$, positive exactly
when $\beta<\tfrac34$ (\path{fourGroupCertPositiveCost_gap_sign_iff}), so
even with positive resource costs the relative welfare of the two specified
paths reverses with the planner's discount factor. Both are existence
certificates: they refute the universal ``region~(R) implies a
positive gap'' and say nothing about magnitudes. Politically weighted payoff
coefficients therefore do not identify the comparison by aggregate welfare;
the region conditions constrain cross-partials and curvature, not the
resource burn against the groups' benefits.

\paragraph{From a feasible comparison to the planner value.}
For a polity, a transition $\F$ and an initial state \(x_0\), let
\(\mathcal P(x_0)\) be the set of
cube-valued path triples \((x_t,a_t,s_t)_{t\ge0}\) starting at \(x_0\) and
satisfying the transition law \(x_{t+1}=\F(x_t,s_t)\), and define
\[
\mathcal W^\star_\beta(x_0)
:=
\sup_{(x,a,s)\in\mathcal P(x_0)}
\sum_{t=0}^{\infty}\beta^tW(x_t,a_t,s_t).
\]
On the cube \(\lvert W\rvert\) is at most a bound in the coefficients
(\path{plannerWBound}; \(13/2\) at the certificate-complete polity, whose
welfare function is \eqref{eq:planner-W-fgc}), so this supremum is finite.
(The bound is a property of the polity's coefficients; it has no counterpart
for the interest-group microfoundations of the family below, whose welfare coefficients
are unbounded across microfoundations.) Write
\(\varphi(a)=(\beta^{\mathrm{base}}_S+c_K-c_J)a-k_Ta^2\) for the compliance part
of $W$, \(a^\ast\) for a maximizer of \(\varphi\) on \([0,1]\), which exists by
compactness (\path{exists_complianceFlow_max}), and
\(Y(x_0)=\sum_t\beta^ty_t\) for the discounted state along the zero-sanction
path from \(x_0\).

\begin{proposition}[The planner's optimum and the exact loss]
\label{prop:planner-optimum-lower-bound}
\begin{enumerate}[label=(\roman*)]
\item For every interest-group polity, every transition
$\F\colon[0,1]^2\to[0,1]$ satisfying \textup{(F2)--(F3)}, every
\(\beta\in[0,1)\) and every \(x_0\in[0,1]\), write $B=b_J-b_K$ and suppose
either (a) $B\ge0$, the transition has an upper bound $L$ on its sanction slope,
$\F(x,s_2)-\F(x,s_1)\le L(s_2-s_1)$ for $s_1\le s_2$, and an upper bound
$\lambda$ on its state slope along the zero-sanction section,
$\F(y_2,0)-\F(y_1,0)\le\lambda(y_2-y_1)$ for $y_1\le y_2$, with
$\beta\lambda<1$, and
\[
g_J-q_K+\frac{\beta BL}{1-\beta\lambda}\le d_s ;
\]
or (b) $B\le0$ and $g_J-q_K\le d_s$. Then the zero-sanction path with the
constant compliance $a^\ast$ attains the planner's optimum,
\[
\mathcal W^\star_\beta(x_0)=B\,Y(x_0)+\frac{\varphi(a^\ast)}{1-\beta},
\]
and the optimum exceeds the value of the corner path of full
sanctions against full non-compliance by
$H/(1-\beta)+B\sum_t\beta^tz_t+\varphi(a^\ast)/(1-\beta)$, with $H$ and $z_t$
as in \Cref{prop:affine-welfare-gap}. Under the partial-adjustment law \eqref{eq:affine},
for either sign of $B$, the condition
$g_J-q_K+\beta\delta B/(1-\beta(1-\delta))\le d_s$ gives
$\mathcal W^\star_\beta(x_0)=Bx_0/(1-\beta(1-\delta))+\varphi(a^\ast)/(1-\beta)$.
\item At the certificate-complete interest-group economy, for every
\(x_0\in[0,1]\) and \(\beta\in[0,1)\), the zero-sanction path attains the
planner's optimum, \(\mathcal W^\star_\beta(x_0)=x_0/(1-\beta/2)\), and
\[
\mathcal W^\star_\beta(x_0)-V^E(x_0;\beta)
=
\frac{23}{10(1-\beta)}+\frac{1}{1-\beta/2}
>0,
\]
where \(V^E\) is the discounted welfare of the play of the constructed
equilibrium.
\end{enumerate}
\end{proposition}

The planner condition of (i) says that the first unit of sanction does not pay socially once the
political shift it causes is counted. A unit of sanction gives the sender's
groups $g_J-q_K$ net and raises next period's state by at most $L$; along the
zero-sanction path that follows, the planner values the higher state at most
$\beta B/(1-\beta\lambda)$ per unit when $B\ge0$ and at most zero when
$B\le0$, since a planner with $B<0$ counts a higher state as a loss. If these
together do not cover the target's direct sanction cost $d_s$, the zero-sanction
path is optimal: the planner gains nothing by sanctioning. The quadratic cost
$k_Ss^2$ and the compliance--sanction cross term only strengthen the conclusion.

\begin{proof}
(i) Let \(v(x)=B\,Y(x)+\varphi(a^\ast)/(1-\beta)\). At every feasible triple,
\(W(x,a,s)+\beta v(\F(x,s))-v(x)\) is the sum of the political term
\((g_J-q_K-d_s)s+\beta B\bigl(Y(\F(x,s))-Y(\F(x,0))\bigr)\) and
\(\varphi(a)-\varphi(a^\ast)-k_Ss^2-w_{\mathrm{pl}}sa\), which is nonpositive by
the choice of \(a^\ast\) and the polity's signs $k_S>0$ and
$w_{\mathrm{pl}}\ge0$ (\path{bellman_le_of_politicalMargin}). Where the
political term is nonpositive at every feasible state and sanction, the
supersolution bound of \path{Foundations/DiscountedChordSum.lean} bounds every
feasible value by \(v(x_0)\), and the zero-sanction path with compliance
\(a^\ast\) attains it (\path{plannerOptimalValue_eq_of_politicalMargin}). For
$B\ge0$, $Y$ is non-decreasing, two zero-sanction paths stay within
$\lambda^t$ times their initial distance, and so
$Y(\F(x,s))-Y(\F(x,0))\le Ls/(1-\beta\lambda)$, which with the condition of (a)
makes the political term nonpositive (\path{zeroStateSum_sub_le},
\path{plannerOptimalValue_eq_of_upperChords}; $\lambda\ge0$ is derived, not
assumed). For $B\le0$ a sanction raises the state, which a planner with
$B\le0$ values at no more than zero, so (b) suffices and no slope bound is used
(\path{plannerOptimalValue_eq_of_nonpos}). Under the partial-adjustment law
$Y(x)=x/(1-\beta(1-\delta))$ exactly (\path{plannerOptimalValue_eq_affineF}).
The loss is the accounting identity of \Cref{prop:affine-welfare-gap} plus the
capitalized compliance flow (\path{plannerOptimalValue_loss_of_upperChords},
\path{plannerOptimalValue_loss_of_nonpos},
\path{plannerOptimalValue_loss_affineF}).

(ii) At the certificate-complete economy $B=1$, $g_J=q_K$ and
\(\varphi\equiv0\), and the planner condition of (i) under the partial-adjustment law at \(\delta=\tfrac12\) reads
\(\beta/(2-\beta)\le\tfrac{14}{5}\), which holds at every \(\beta\in[0,1)\);
so (i) gives the optimum (\path{plannerOptimalValue_eq_affineF},
\path{fgcPlannerOptimalValue_eq_general}). The verification at the example economy is
direct.
Let \(v_0(x)=x/(1-\beta/2)\). At every feasible triple,
\(W(x,a,s)+\beta v_0(\F(x,s))-v_0(x)
=s\bigl(\tfrac{\beta}{2-\beta}-\tfrac{14}{5}\bigr)-\tfrac{s^2}{2}-\tfrac{11}{5}sa\le0\),
with equality at \(s=0\) (\path{fgc_planner_bellman_upper}). Multiplying by
\(\beta^t\) and summing over \(t<N\) along any feasible path telescopes to
\(\sum_{t<N}\beta^tW_t+\beta^Nv_0(x_N)\le v_0(x_0)\); the terminal term
vanishes because the state stays in \([0,1]\), and the series converges by
the cube bound, so every feasible value is at most \(v_0(x_0)\)
(\path{fgc_plannerPath_value_le}, through the generic verification lemma of
\path{Foundations/InfiniteHorizonVerification}). The zero-sanction path
attains \(v_0(x_0)\) by \eqref{eq:planner-values}, so the supremum equals it
(\path{fgcPlannerOptimalValue_eq}), and the exact difference in
\eqref{eq:planner-values} gives the loss of equilibrium play relative to the
specified feasible zero-sanction path
(\path{fgc_plannerOptimalValue_gap_eq_equilibrium}).
\end{proof}

At the certificate-complete economy the optimum is attained by the
zero-sanction path, but not uniquely: at
\(s=0\) every compliance choice ties, and at \(\beta=0\) future controls are
not ranked.

The planner condition is sufficient, and it fails at the hawk-scaled and
positive-cost polities above. At the first, $g_J-q_K=4>\tfrac{14}5=d_s$, so it
fails at every discount; at the second, where $d_s=\tfrac1{10}$, the condition under the
partial-adjustment law reads $\beta/(2-\beta)\le\tfrac1{10}$ and fails for
$\beta>\tfrac2{11}$. There sanctions can pay socially: at the first polity at
every discount, and at the second at $\beta=\tfrac9{10}$, the corner path of full sanctions against full
non-compliance is worth strictly more to the planner than the zero-sanction path at
full non-compliance (\path{fourGroupCertHawkScaled_gap_neg},
\path{fourGroupCertPositiveCost_gap_at_nine_tenths}). At both polities
$\beta^{\mathrm{base}}_S=0$, $c_J=c_K$ and $k_T=0$, so the compliance part
$\varphi$ of $W$ is identically zero and every zero-sanction path is worth the
same. Since the corner path is feasible, an immediate consequence is that no
zero-sanction path attains the planner's optimum there; this step is an
inference from the two comparisons, not a separate machine-checked result.

Evaluating the exact gap at $\beta_S=\tfrac12$ and at
$\beta_T=\tfrac9{10}$ returns the two capitalized numbers reported
there. The planner condition in general form is
\Cref{prop:zero-sanction-planner}, and the sensitivity of the verdict to how $W$
is defined is \Cref{prop:welfare-rescaling}; the verdict itself --- a
lock-in that is politically rational and socially wasteful --- is the
sanctions-persistence analogue of the inefficient-persistence verdicts in
the dynamic political-economy tradition
\citep{AcemogluRobinson2008,BattagliniCoate2008,HasslerEtAl2003}.

\paragraph{Interest-group microfoundations of the quadratic--linear family under the lifting condition.}
Everything above, the cube bound $13/2$ included, concerns affine interest-group
polities that sit next to one example economy. The remainder of this
section leaves them. The same non-identification holds on the whole family of
\Cref{sec:four-regimes}, through a separate accounting object that leaves the
interest-group polity of this appendix untouched (that polity has no direct
state benefit $wx$ and requires strict signs on the state--sanction cross terms, while the family allows
$r=0$ and $n=0$). An \emph{interest-group microfoundation} of a family member's government payoffs consists of
target groups with state valuations $b_J,b_K>0$ and a compliance transfer
$\kappa>0$ between them, baseline political weights $t_J,t_K>0$ with a common
sanction slope $\nu_T\ge0$, and sender groups gaining $g_H>0$ and losing $Q>0$
per unit of sanction, with weights $(u_0+Q+rx)/g_H$ and one, subject to
\[
t_Jb_J-t_Kb_K=v,\qquad \nu_T(b_J-b_K)=n,\qquad \kappa(t_J-t_K)=\ell .
\]
The direct state benefit $wx$, the sanctioning cost $\tfrac q2s^2$, the
compliance--sanction cross term $kas$ and the target's sanction cost $cs$ are kept at unit
weight. Every weight is strictly positive on the feasible cube.

\begin{lemma}[Interest-group microfoundations: representation and existence]
\label{lem:welfare-completion}
For every interest-group microfoundation the politically weighted group sums equal
$\Us(x,a,s)$ and $\Ut(x,a,s)$ at every state and action, so the government
game and every constructed equilibrium are the member's own, and aggregate
welfare is
\[
W(x,a,s)=\bigl(w+b_J-b_K\bigr)x+\bigl(g_H-Q-c\bigr)s-\tfrac q2s^2-kas .
\]
Every member has a microfoundation with any prescribed net sanction coefficient
$\eta=g_H-Q-c$ and any $b_J-b_K>0$. Microfoundations with $b_J=b_K$ exist exactly
at the members with $n=0$, and there at every $\eta$.
\end{lemma}

\begin{proof}
The representation is algebra in the three displayed equations
(\path{govS_eq}, \path{govT_eq}, \path{welfare_eq}); existence takes $b_K=1$,
$b_J=1+B$, $\kappa=2\ell b_J/v$, $t_K=v/(2B)$, $t_J=t_K+\ell/\kappa$,
$\nu_T=n/B$, $Q=|c+\eta|+1$ and $g_H=Q+c+\eta$
(\path{exists_welfareCompletion}); these are choices inside an existence
proof, not restrictions on the class, and the class does not require
$b_J>b_K$. Since $\nu_T(b_J-b_K)=n$, equal valuations force $n=0$; and at
$n=0$ the record $b_J=b_K=1$, $t_K=1$, $t_J=1+v$, $\kappa=\ell/v$, $\nu_T=0$,
with the same $Q$ and $g_H$, is a microfoundation at every $\eta$, whatever the
transition and the state--sanction cross term $r$ (\path{exists_completion_bJ_eq_bK_iff},
\path{exists_equalValuation_completion}). The microfoundation of
\Cref{cor:planner-efficient-boundary} is of this kind.
\end{proof}

The family has its own version of \Cref{prop:welfare-nonidentification}, with
these microfoundations in place of the accounting representations used there.

\begin{proposition}[Welfare non-identification on the family]
\label{prop:family-welfare-nonidentification}
Take a member of the quadratic--linear family in the class that satisfies
the lifting condition, the entrenchment condition and the forward-looking relief condition
\textup{(RL)}, and at which the sender's response at full non-compliance at its
constructed value is non-decreasing in the state, as under the increasing-differences condition \eqref{eq:rally-margin} (\Cref{app:proofs:placable}).
There are two interest-group microfoundations such that, at every initial state
and every common planner discount factor $\beta\in[0,1)$, discounted aggregate
welfare along the four constructed equilibria is ranked
\[
11\succ10\succ00\sim01\quad\text{by one},\qquad
00\sim01\succ10\succ11\quad\text{by the other}.
\]
\end{proposition}

\begin{proof}[Proof of \Cref{prop:family-welfare-nonidentification}]
For the two ranking microfoundations take $b_J-b_K=1$, so that the state
coefficient is $K=w+1>0$. The constructed equilibria of cases $00$ and $01$ have
the zero sanction and one state path, and $W(x,a,0)$ does not depend on $a$,
so they tie. Along the other two the target maintains non-compliance, the target-only
sanction is the stage vertex, at least $u_0/q$, and the forward-looking
sanction exceeds it by at least the foresight premium \eqref{eq:sender-wedge} at every
date, while the state paths are ordered $00\le10\le11$. With $\eta=q+1$ every
flow comparison is strictly positive in the direction of more sanctions; with
$\eta$ below $-(w+1)/\min\{u_0/q,\text{premium}\}-1$ it is strictly negative,
since a state difference is at most one. The flows are bounded, so the strict
termwise comparison survives discounting, including at $\beta=0$
(\path{exists_completions_reversing_regime_ranking_of_response_monotone}; the
form under the increasing-differences condition, \path{exists_completions_reversing_regime_ranking}, is its
corollary).
\end{proof}

\begin{corollary}[A planner-efficient boundary]
\label{cor:planner-efficient-boundary}
Let a member have $r=n=0$. Then there is a microfoundation with $W=\Us$ such that,
at the planner discount $\beta_S$, the play of every equilibrium in $\mathcal E$
at which the target maintains full non-compliance --- full non-compliance and its sanction rule,
at every state --- attains the planner's optimum against every
time-dependent plan over both actions, from every start.
\end{corollary}

\begin{proof}
Take
$b_J=b_K=1$, $t_J=1+v$, $t_K=1$, $\kappa=\ell/v$, $Q=1$ and $g_H=1+c+u_0$;
then $W=\Us$, the sender's constructed value satisfies
$W(x,a,s)+\beta_S V_S(\F(x,s))\le\Us(x,0,s)+\beta_S V_S(\F(x,s))\le V_S(x)$
with equality at full non-compliance and the dynamic best response to it, and the bounded
verification theorem of \path{Foundations/InfiniteHorizonVerification.lean}
concludes (\path{exists_planner_efficient_completion}). That inequality holds
on all of $r=n=0$. Every full-non-compliance equilibrium has the constructed values
and plays the constructed response, so the compared plan is that
equilibrium's policy and sanction rule, with neither the entrenchment condition nor the increasing-differences condition
(\path{exists_planner_efficient_completion_of_policy_zero}); the form at the
constructed equilibrium under those two conditions is
\path{exists_planner_efficient_completion_at_equilibrium}, the flatness the
increasing-differences condition forces at $r=0$ is \path{flat_of_margin_of_r_eq_zero}, and the
inhabitant is \path{placableBoundary_planner_efficient_at_equilibrium}. The
statement concerns one microfoundation and one planner discount: it is an
existence boundary for any universal strict-inefficiency sentence, not
generic efficiency.
\end{proof}

Such an equilibrium exists exactly when full non-compliance is a target best reply
at the constructed values; under the entrenchment condition and a non-decreasing
response the constructed equilibrium of \Cref{prop:placable-family} is one.
The corollary needs neither the entrenchment condition nor the increasing-differences condition. Through the increasing-differences condition it reaches only flat
transitions: at $r=0$ that condition forces a flat transition certificate,
$M_{\mathrm{mix}}=0$ and $M_{\mathrm{cx}}\,g_{\mathrm{hi}}=0$, as under the
partial-adjustment law and every additively separable transition with convex state
sections. The boundary is inhabited by the member with $r=n=0$ whose transition is the partial-adjustment law
(\Cref{app:formalization}); no member with $r=n=0$, a non-flat transition and a
full-non-compliance equilibrium is exhibited.

The quantifier order is: for every member there exist two microfoundations that
work for every start and every planner discount. The microfoundations depend on
the member; no claim is made about every microfoundation, about all equilibria, or
about magnitudes. Three accounting choices carry the statement and are choices,
not consequences of the government game: $wx$ counted at unit social weight;
compliance as a pure transfer between the two target groups, which is what
makes the constructed equilibria of cases $00$ and $01$ tie; and group benefits unrestricted in size, with
the inverse political weight compensating. The microfoundations retain the
compliance--sanction cross term $kas$ at unit weight, while the compliance transfer $\kappa a$
cancels between the two target groups; counting a term at unit weight is a
convention, not a demonstration that it is a resource loss. A bounded class
is treated in \Cref{prop:welfare-interval} below. A ranking
of the constructed equilibria is not a planner-efficiency verdict: under the first microfoundation a
planner would want the maximal sanction, not the interior equilibrium one, so
the forward-looking equilibrium is still planner-inefficient there;
\Cref{cor:planner-efficient-boundary} is what shows that inefficiency itself
is not universal.

\paragraph{What a restriction on the microfoundations identifies.}
With cardinal group benefits free and political weights compensating,
non-identification is what one should expect. The next proposition says what a
restriction buys. Along the constructed equilibria of the two non-compliance cases compliance is zero, so the
welfare gap is affine in the net sanction coefficient; the restriction is on the
sender side only, and the target side of the microfoundation is held fixed.

\begin{proposition}[Partial identification with fixed exposure and bounded hawk weights]
\label{prop:welfare-interval}
Let a member satisfy the entrenchment condition, with the sender's response at
full non-compliance at its constructed value non-decreasing in the state (as under
the increasing-differences condition \eqref{eq:rally-margin}; \Cref{app:proofs:placable}), fix an
initial state and a planner discount
$\beta\in[0,1)$, and let $\Delta X$, $\Delta S$ and
$\Delta S_2$ be the discounted differences between the constructed equilibria of the forward-looking and the
target-only cases in the state, the sanction and the squared sanction.
\begin{enumerate}[label=(\roman*)]
\item For every interest-group microfoundation,
$V^{11}-V^{10}=K\,\Delta X+\eta\,\Delta S-\tfrac q2\Delta S_2$ with
$K=w+b_J-b_K$ and $\eta=g_H-Q-c$; the three differences do not depend on the
microfoundation, and $\Delta S\ge m/(1-\beta)>0$ with $m$ the foresight premium
\eqref{eq:sender-wedge}.
\item Fix the commercial lobby's exposure $Q=\bar Q>0$ and require the hawks'
political weight to satisfy $\underline\theta\le(u_0+\bar Q+rx)/g_H\le\bar\theta$
at every feasible state, with $0<\underline\theta$, $0<\bar\theta$. This class is
non-empty if and only if
$\underline\theta\,(u_0+\bar Q+r)\le\bar\theta\,(u_0+\bar Q)$, and on it $\eta$
ranges over exactly
$\bigl[(u_0+\bar Q+r)/\bar\theta-\bar Q-c,\;(u_0+\bar Q)/\underline\theta-\bar Q-c\bigr]$,
every value being attained with the target side, hence $K$, unchanged.
\item With the target side fixed, the set of welfare gaps over the class is the
closed interval with endpoints
$K\Delta X-\tfrac q2\Delta S_2+\eta\,\Delta S$ at the two ends of that range.
On a non-empty class, among its microfoundations with the given state coefficient
$K$ --- a set that contains those with the given target side --- the
equilibrium of the forward-looking case ranks strictly first at every microfoundation exactly when
the lower endpoint is positive, and strictly last exactly when the upper
endpoint is negative; both strict rankings occur exactly when the endpoints have
opposite strict signs, and then at microfoundations that keep the whole target side.
An endpoint at zero leaves a weak ordering, and the interval $\{0\}$ a tie.
\item A weight band alone need not identify the sign. For any band straddling
one, $\underline\theta<1<\bar\theta$, unrestricted exposure permits both
rankings: with the target side kept, one microfoundation with weights inside the
band ranks the equilibrium of the forward-looking case strictly above that of the target-only case and
another ranks it strictly below. Nothing is claimed for a band that does not
straddle one.
\end{enumerate}
\end{proposition}

\begin{proof}
(i) is \Cref{lem:welfare-completion} along paths on which $as=0$, with bounded
flows summed termwise, and the foresight premium at every date
(\path{regimeValue11_sub_10_eq}, \path{dS_ge}). (ii): the weight rises with the
state because $r\ge0$, so the two bounds are
$(u_0+\bar Q+r)/\bar\theta\le g_H\le(u_0+\bar Q)/\underline\theta$; the
representation equations involve neither $g_H$ nor $Q$, so any $g_H$ in that
range can be installed without touching the target side
(\path{hawkBounded_iff}, \path{hawkBounded_nonempty_iff},
\path{hawkBounded_eta_mem}, \path{exists_withSender_hawkBounded_eta}, whose
attaining microfoundation is the given one with $g_H$ and $Q$ replaced and every
other field kept). (iii) follows from
(i), (ii) and $\Delta S>0$ (\path{gap_mem_interval}, \path{gap_attained_withSender},
\path{forall_gap_pos_iff}, \path{forall_gap_neg_iff}, \path{exists_both_signs_iff},
and \path{exists_withSender_both_signs_of_straddle}, which names the two
microfoundations). (iv): with $g_H=(u_0+Q)/\underline\theta$
the weights stay in the band once $Q$ is large and
$\eta=(1/\underline\theta-1)Q+u_0/\underline\theta-c\to+\infty$; with
$g_H=(u_0+Q+r)/\bar\theta$ they stay in the band and $\eta\to-\infty$
(\path{exists_withSender_weightBounded_eta_gt},
\path{exists_withSender_weightBounded_eta_lt},
\path{exists_withSender_weightBounded_reversal}; each names its microfoundation as
the given one with $g_H$ and $Q$ replaced). Each of the theorems on the constructed equilibria
cited is stated under the increasing-differences condition and is the corollary of a sibling that
assumes only the non-decreasing response
(\path{regimeValue11_sub_10_eq_of_response_monotone},
\path{dS_ge_of_response_monotone}, \path{gap_mem_interval_of_response_monotone},
\path{gap_attained_withSender_of_response_monotone},
\path{forall_gap_pos_iff_of_response_monotone},
\path{forall_gap_neg_iff_of_response_monotone},
\path{exists_both_signs_iff_of_response_monotone},
\path{exists_withSender_both_signs_of_straddle_of_response_monotone},
\path{exists_withSender_weightBounded_reversal_of_response_monotone}); the
theorems on the microfoundations alone read neither the entrenchment condition nor the increasing-differences condition.
\end{proof}

The restricted class is a definition over the microfoundation's own data, not a
hypothesis on the governments' game: payoffs and constructed equilibria are
untouched. It is one restriction among many; holding the target side fixed is
substantive, and the proposition identifies the ranking of two constructed equilibria, not a
planner optimum. The interval is conditional on accounting information the
governments' behavior does not supply; its sharpness does not supply it
either. What it shows is where the information has to come from:
group exposures, not political weights alone.
